\section{Implications for Individuals, Organizations, and Society}
\label{sec:implications}

A measurement is only worth the decisions it improves. This section sets out what follows from
the accounts for the people the exposure reaches, for the organizations that have to govern it,
and for the wider public interest. We separate what the evidence supports from what it does not,
because recommendations are the easiest part of a paper of this kind to overreach.

\paragraph{For the people the exposure reaches.} The household that services a mortgage, a car
loan and a student loan out of one wage is not an edge case in this system; it is the system's
modal unit. Our measurements place it there. One practical consequence is that liquid assets
relative to income are not monotonic in pay: measured that way they are lowest one quintile above
the bottom, among households carrying the commitments of those above them without the income to
service them. We stop short of calling those households the most vulnerable, because that
depends on expenditure and commitments the accounts do not carry, but a relief design that
assumes liquidity rises with pay is assuming a pattern the data do not show. Relief designed
around occupation faces a separate problem: the accounts show the exposure is spread almost
evenly across borrowers rather than concentrated in identifiable jobs.

There is a second consequence, and it runs the other way. Because the claims are spread more
evenly than the wages, the exposure is not concentrated in a group that could be identified and
advised. We are careful about what follows. The accounts measure a stock and its attribution;
they do not survey the instruments available to a household, so they cannot establish that no
such instrument exists. What they do establish is narrower: the exposure is common rather than
sorted, so any remedy that works by identifying the exposed will have little to sort on. Whether
public instruments or disclosure are the right answer is a question about costs and
counterfactuals that these accounts do not price.

\paragraph{For managers and organizations.} The governance question these accounts pose is
narrow and answerable: can an institution state, in a form a board can read, how much of its
income and collateral depends on wages continuing to be paid? Most cannot, because no standing
statistic reports it. That is a measurement gap an institution can close for itself from its own
book using the coefficients in Table~\ref{tab:bridge}, and the exercise is cheap.

What the evidence does \emph{not} support is tightening household underwriting in response. Three
measured grounds stand against it: the debt service result is partly mechanical, since
underwriting already caps debt service to income; the exposure is spread almost evenly across
borrowers, so screening on it buys a lender little; and the one pre-registered institutional
test we ran was a null, which Section~\ref{sec:limitations} reports in full. That last point
is a null on one outcome in one window, not a demonstration that the accounts can never carry
institution-level predictive content, and we do not read it as the stronger claim. Wage
dependence is a concentration to be governed and disclosed, not a credit category to be
underwritten.

\paragraph{Strategic implications.} The holder map makes an ownership question explicit that is
usually left implicit. If every holder class carries wage-attributed claims in some proportion,
then hedging inside the system requires a counterparty whose own exposure is lower, and the map
says which classes those would be rather than that none exist. We do not claim the exposure
cannot be diversified away, and we do not rank instruments: the accounts carry no prices, no
contract terms and no costs, so a statement that one instrument has the best ratio of value to
cost is not available from them. What the map supports is weaker and still useful: an
institution can read its own wage-attributed share off a published construction, which makes the
number auditable rather than discretionary.

The second strategic point concerns timing. The federal share of the wage-backed stock has
more than doubled from its trough, \result{SovUnionTrough} percent in \result{SovUnionTroughYear} against
\result{SovereignUnion} percent in 2025, a factor of \result{SovUnionRiseFromTrough}. Across the whole series it
has not doubled: it was \result{SovUnion1947} percent in 1947, so the series is a U and the rise is
from the middle of it rather than from the start. The whole net increase since 2008 sits in the
debtor position rather than in new guarantees. A position acquired that way is acquired without a
decision having been taken about it. Making it visible is the precondition for deciding whether
to keep it.

\paragraph{For society and equity.} The claims that wages pay are spread far more evenly than the
wages themselves. The implication is a balance-sheet one and we keep it there: a proportional
fall in pay lands against claim stocks that are larger relative to wages further down the
distribution, so the ratio of claims to wages rises as pay falls. That is not the same as showing
the fall is regressive in welfare, in default or even in debt service, each of which depends on
payment schedules, maturity, transfers, other assets and expenditure, none of which these
accounts carry. What the structure of the claim stock settles is the gradient of the stock. The
incidence of a shock on wellbeing or on arrears is a further question and we do not answer it
here.

Two refusals belong here as well, because a statistic that locates wage risk could be used
against the people it describes. These accounts are built at the level of claim classes and
holders, not individuals, and nothing in them supports pricing a person's credit by what they do
for a living. Occupation is not a prohibited basis under the Equal Credit Opportunity Act, whose
list of prohibited bases does not include it, but because occupation correlates with protected
classes, occupation-based underwriting in housing credit exposes a lender to a
discriminatory-effects claim under the Fair Housing Act, on which the lender must carry a legally
sufficient justification. Combined with the finding that the exposure is spread almost evenly
across borrowers, the risk-adjusted case for it is weak and the case for portfolio-level
disclosure is strong. A statistic that locates wage risk should size public instruments. It
should not price individuals out of credit for what they do for a living.
